\documentclass[11pt,a4paper]{article}
\usepackage[utf8]{inputenc}
\usepackage{amsmath,amssymb,amsthm}
\usepackage{geometry}
\usepackage{booktabs}
\usepackage{hyperref}
\usepackage{algpseudocode}

\geometry{margin=2.5cm}

\newtheorem{theorem}{Theorem}
\newtheorem{lemma}[theorem]{Lemma}

\title{Problem 403: Lattice Points Enclosed by Parabola and Line}
\author{}
\date{}

\begin{document}
\maketitle

\section{Problem Statement}

For integers $a$ and $b$, define
\[
D(a,b) = \bigl\{(x,y) \in \mathbb{R}^2 \mid x^2 \le y \le ax + b \bigr\}.
\]
Let $L(a,b)$ denote the number of lattice points in $D(a,b)$.  Define
\[
S(N) = \sum_{\substack{|a| \le N,\; |b| \le N \\ \operatorname{Area}(D(a,b)) \in \mathbb{Q}}} L(a,b).
\]
Given $S(5) = 344$ and $S(100) = 26\,709\,528$.

\medskip\noindent\textbf{Find} $S(10^{12}) \bmod 10^8$.

\section{Mathematical Foundation}

\begin{theorem}[Rational area condition]
\label{thm:rational}
The area of $D(a,b)$ is rational if and only if $\Delta = a^2 + 4b$ is a non-negative perfect square.
\end{theorem}

\begin{proof}
The intersections of $y = x^2$ and $y = ax+b$ occur at roots of $x^2 - ax - b = 0$, with $r_{1,2} = (a \mp \sqrt{\Delta})/2$ where $\Delta = a^2 + 4b$.  The enclosed area is
\[
\operatorname{Area} = \int_{r_1}^{r_2}(ax+b-x^2)\,dx = \frac{(r_2-r_1)^3}{6} = \frac{\Delta^{3/2}}{6}.
\]
Since $\Delta^{3/2} = \Delta \cdot \sqrt{\Delta}$, this is rational iff $\sqrt{\Delta} \in \mathbb{Q}$, iff $\Delta$ is a perfect square (or zero).
\end{proof}

\begin{lemma}[Parity constraint]
If $\Delta = s^2$ and $b = (s^2 - a^2)/4 \in \mathbb{Z}$, then $a \equiv s \pmod{2}$.
\end{lemma}

\begin{proof}
$4 \mid (s-a)(s+a)$ requires $s \equiv a \pmod{2}$, since if they have different parity then $(s-a)(s+a)$ is odd.
\end{proof}

\begin{theorem}[Closed form for $L$]
\label{thm:L}
When $\Delta = s^2$ with $s \ge 0$,
\[
L(s) = \frac{(s+1)(s^2 - s + 6)}{6}.
\]
\end{theorem}

\begin{proof}
The roots $r_1 = (a-s)/2$ and $r_2 = (a+s)/2$ are integers.  Substituting $t = x - r_1$:
\[
L(a,b) = \sum_{t=0}^{s}\bigl[t(s-t)+1\bigr] = \frac{s(s+1)(s-1)}{6} + (s+1) = \frac{(s+1)(s^2-s+6)}{6}. \qedhere
\]
\end{proof}

\begin{lemma}[Reparametrization]
Setting $a = s + 2u$, we get $b = -u(s+u)$ with constraints $|s+2u| \le N$, $|u(s+u)| \le N$, $s \ge 0$.
\end{lemma}

\begin{proof}
From $s^2 = a^2 + 4b$: $b = (s^2-(s+2u)^2)/4 = -u(s+u)$.
\end{proof}

\begin{theorem}[Prefix sum]
\label{thm:prefix}
$P(K) := \sum_{s=0}^{K} L(s) = \frac{(K+1)(K+2)(K^2 - K + 12)}{24}$.
\end{theorem}

\begin{proof}
Sum $L(s) = (s+1)(s^2-s+6)/6$ using standard power-sum identities for $\sum(s+1)$, $\sum s(s+1)$, $\sum s^2(s+1)$, then factor the result.
\end{proof}

\begin{theorem}[Sub-linear computation]
$S(N)$ is computable in $O(\sqrt{N})$ time by splitting the sum over $u$:
\begin{enumerate}
  \item $u=0$: contribution $P(N)$.
  \item $u \ge 1$: direct iteration for $u \le \lfloor\sqrt{N}\rfloor$ (since $u^2 \le N$).
  \item $u \le -1$, $v=-u$: for $v \le \sqrt{N}$, direct; for $v > \sqrt{N}$, group by $q = \lfloor N/v \rfloor$, giving $O(\sqrt{N})$ groups with $O(1)$ work each via Faulhaber's formulas.
\end{enumerate}
\end{theorem}

\begin{proof}
For $u \ge 1$: $|u(s+u)| \le N$ with $s \ge 0$ forces $u^2 \le N$.  For $v > \sqrt{N}$: $\lfloor N/v\rfloor$ takes $O(\sqrt{N})$ distinct values.  Within each group, contributions are degree-4 polynomials in~$v$, summable via $\sum v^k$ ($k = 0,\ldots,4$) in $O(1)$.
\end{proof}

\section{Algorithm}

\begin{algorithmic}[1]
\Function{$S_{\mathrm{mod}}$}{$N, M$}
  \State $\mathit{result} \gets P(N) \bmod M$ \Comment{Case $u=0$}
  \For{$u = 1$ \textbf{to} $\lfloor\sqrt{N}\rfloor$} \Comment{Case $u \ge 1$}
    \State $s_{\max} \gets \min\!\bigl(\lfloor(N-u^2)/u\rfloor,\; N-2u\bigr)$
    \If{$s_{\max} \ge 0$} $\mathit{result} \mathrel{+}= P(s_{\max})$ \EndIf
  \EndFor
  \For{$v = 1$ \textbf{to} $\lfloor\sqrt{N}\rfloor$} \Comment{Case $v \le \sqrt{N}$}
    \State Compute valid $s$-range $[s_{\mathrm{lo}}, s_{\mathrm{hi}}]$; add $P(s_{\mathrm{hi}}) - P(s_{\mathrm{lo}}-1)$
  \EndFor
  \For{$q = 1$ \textbf{to} $\lfloor\sqrt{N}\rfloor$} \Comment{Case $v > \sqrt{N}$, grouped by $q$}
    \State Sum degree-4 polynomial in $v$ via Faulhaber formulas over group
  \EndFor
  \State \Return $\mathit{result} \bmod M$
\EndFunction
\end{algorithmic}

\section{Complexity Analysis}

\begin{itemize}
  \item \textbf{Time:} $O(\sqrt{N})$.  All cases iterate over at most $O(\sqrt{N})$ values, each with $O(1)$ work.
  \item \textbf{Space:} $O(1)$.
\end{itemize}

\section{Answer}

\[
\boxed{18224771}
\]

\end{document}
